\section*{РОЗДІЛ 1. ТЕОРЕТИКО-МЕТОДОЛОГІЧНІ ЗАСАДИ ДОСЛІДЖЕННЯ}

\subsection*{1.1. Концептуальний апарат: реактивне та проактивне управління}

Теоретичним фундаментом дослідження є дихотомія реактивного та проактивного підходів до управління, яка набуває особливого значення в контексті цифрової трансформації публічного сектору.

\textbf{Реактивне управління} (reactive management) характеризується як модель організаційної поведінки, при якій управлінські рішення та дії ініціюються у відповідь на вже виниклі проблеми, зовнішні запити або настання дедлайнів. За визначенням П. Данліві та співавторів, реактивні адміністративні системи функціонують переважно як «буфери обміну інформацією», активуючись лише під впливом зовнішнього стимулу \cite{Dunleavy2006DigitalEraGovernance}. У контексті адміністрування науки це проявляється у формуванні звітів post factum, ручній компіляції даних за запитом керівництва або контролюючих органів, відсутності системного моніторингу.

Дослідження OCLC Research, проведене Р. Брайантом та колегами, виявило системні характеристики реактивних моделей управління дослідницькою інформацією в університетах США \cite{Bryant2017RIMDefining}:

\begin{itemize}[leftmargin=*,itemsep=0.2em]
\item \textbf{Децентралізація без координації:} управління науковою інформацією відбувається на рівні окремих факультетів або кафедр без єдиної точки контролю та агрегації;
\item \textbf{Інформаційні силоси:} дані про публікації, гранти, кадрові ресурси зберігаються в розрізнених системах, що унеможливлює комплексний аналіз;
\item \textbf{Ручна компіляція CV:} оновлення профілів дослідників вимагає індивідуального збору інформації, що призводить до неактуальності даних;
\item \textbf{Фокус на compliance:} діяльність адміністрації орієнтована на формальне виконання звітних вимог, а не на стратегічний розвиток.
\end{itemize}

Ключовою вадою реактивного підходу є проблема «пізнього втручання» (late intervention). Аналогічно до ситуації в академічному консультуванні студентів, де підтримка часто надається лише після загострення проблем до критичного рівня, наукові відділи університетів дізнаються про потреби дослідників занадто пізно --- наприклад, інформація про грантовий конкурс може надходити за кілька днів до дедлайну, унеможливлюючи якісну підготовку заявки.

\textbf{Проактивне управління} (proactive management) представляє альтернативну парадигму, спрямовану на антиципацію потреб та усунення перешкод до їх виникнення. Концепція набула системного обґрунтування в праці \textit{Digital Era Governance} \cite{Dunleavy2006DigitalEraGovernance}, де автори пов'язують її з можливостями цифрових технологій для прогнозування та превентивного реагування.

Ключові принципи проактивного управління:

\begin{enumerate}[leftmargin=*,itemsep=0.2em]
\item \textbf{Антиципація потреб (Anticipating Needs):} використання аналітики даних для передбачення потреб стейкхолдерів. У науковому адмініструванні це означає автоматичне сповіщення дослідників про релевантні грантові можливості, журнали для публікації, конференції --- ще до того, як вони самі почнуть пошук.
\item \textbf{Системний моніторинг:} замість епізодичних перевірок запроваджується постійний моніторинг ключових показників з автоматичним виявленням аномалій та ризиків (наприклад, затримка у виконанні етапів гранту, зниження публікаційної активності).
\item \textbf{Холістичний підхід:} врахування всіх аспектів діяльності дослідника --- від наукових інтересів та публікаційної історії до адміністративного навантаження та кар'єрних цілей. Це вимагає інтеграції даних з різних джерел для формування повного профілю.
\item \textbf{Data-driven decision making:} управлінські рішення базуються на об'єктивних даних, візуалізаціях, предиктивних моделях, а не лише на інтуїції чи традиції.
\end{enumerate}

Емпіричні свідчення ефективності проактивного підходу надає досвід Georgia State University, яка досягла збільшення показників завершення навчання (student retention) на 7 процентних пунктів через впровадження предиктивних інтервенцій на основі ML-моделей \cite{GeorgiaState2019StudentSuccessPredictive}. Система раннього попередження ідентифікувала студентів групи ризику та ініціювала проактивне консультування до того, як проблеми стали критичними.

Трансформація від реактивного до проактивного підходу в адмініструванні науки потребує трьох взаємопов'язаних компонентів:

\begin{longtable}{@{}p{0.22\textwidth}p{0.36\textwidth}p{0.36\textwidth}@{}}
\toprule
\textbf{Компонент} & \textbf{Реактивна модель} & \textbf{Проактивна модель} \\
\midrule
\textbf{Організаційні зміни} & Ізольовані підрозділи & Крос-функціональні команди (науковий відділ + бібліотека + ІТ) \\
\textbf{Технологічна інфраструктура} & Розрізнені бази даних & Інтегровані CRIS-системи з підтримкою ШІ \\
\textbf{Культура управління} & Compliance-орієнтованість & Стратегічне мислення, data-driven culture \\
\bottomrule
\end{longtable}

Framework Jisc для цифрової трансформації вищої освіти \cite{Jisc2023DigitalTransformationFramework} підкреслює, що успішна трансформація потребує ефективного цифрового лідерства на рівні ректорату, належних інвестицій у технології та компетенції, систематичного залучення стейкхолдерів та розвитку цифрової грамотності персоналу.

\subsection*{1.2. Генеративний ШІ як технологічний enabler проактивності}

Реалізація проактивного управління в масштабах університету була практично неможливою до появи технологій штучного інтелекту. Традиційні інформаційні системи забезпечували зберігання та базову аналітику, але не могли обробляти неструктуровані дані (тексти публікацій, описи проєктів), генерувати нові інсайти чи автономно виконувати послідовності дій.

\textbf{Генеративний штучний інтелект} (Generative AI, GenAI) --- це клас моделей машинного навчання, здатних створювати новий контент (текст, код, зображення, структуровані дані) на основі навчання на великих масивах даних. Сучасні великі мовні моделі (Large Language Models, LLM), такі як GPT-4, Claude, Gemini, демонструють здатність до складного міркування, аналізу контексту, генерації структурованих відповідей та, що критично важливо, до використання зовнішніх інструментів \cite{Wang2024SurveyLLMAgents}.

Еволюція можливостей ШІ може бути представлена як перехід через три покоління:

\textbf{Перше покоління --- Rule-based systems:} жорстко запрограмовані правила, неможливість адаптації до нових ситуацій, обмежене застосування для рутинних операцій.

\textbf{Друге покоління --- Predictive AI:} статистичні та ML-моделі для класифікації, прогнозування, виявлення патернів у числових даних. Обмеження: нездатність працювати з неструктурованим текстом, генерувати новий контент.

\textbf{Третє покоління --- Generative AI:} здатність до розуміння та генерації природної мови, роботи з контекстом, виконання складних когнітивних завдань, використання зовнішніх інструментів. LLM стають «reasoning engines», здатними до планування та виконання багатокрокових завдань.

Критичним для нашого дослідження є концепція \textbf{автономних агентів ШІ} (AI Agents). Як систематизовано в оглядовій статті Wang et al.~\cite{Wang2024SurveyLLMAgents}, агенти ШІ принципово відрізняються від «чат-ботів» першого покоління здатністю не лише генерувати відповіді, а й автономно виконувати дії у зовнішньому середовищі:

\begin{itemize}[leftmargin=*,itemsep=0.2em]
\item здійснювати API-запити до баз даних та веб-сервісів;
\item обробляти отримані дані та приймати рішення про наступні кроки;
\item формувати звіти, заповнювати документи, відправляти повідомлення;
\item ітеративно уточнювати результати на основі зворотного зв'язку.
\end{itemize}

Xi et al.~\cite{Xi2023RisePotentialAgents} виділяють три ключові компоненти архітектури ШІ-агента: \textbf{Brain} (LLM як центральний «мозок» для міркування та прийняття рішень), \textbf{Perception} (модулі сприйняття вхідних даних різних модальностей) та \textbf{Action} (модулі виконання дій через зовнішні інструменти).

Механізм \textbf{Function Calling} (виклик функцій) дозволяє LLM автономно вирішувати, коли і які зовнішні інструменти використовувати. Модель отримує декларативний опис доступних функцій (tools) з їх параметрами та, аналізуючи запит користувача, самостійно формує виклики відповідних функцій. Це принципово відрізняється від традиційного програмування, де послідовність дій жорстко задана розробником.

\textbf{Протокол Model Context Protocol (MCP)}, представлений компанією Anthropic у листопаді 2024 року, стандартизує взаємодію LLM-застосунків із зовнішніми джерелами даних та інструментами \cite{Anthropic2024MCP}. У грудні 2025 року протокол передано до Agentic AI Foundation під егідою Linux Foundation з підтримкою провідних технологічних компаній (Google, Microsoft, AWS, Cloudflare).

MCP вирішує критичну проблему $N\times M$ інтеграцій: замість створення окремих конекторів для кожної пари «LLM --- зовнішній сервіс» протокол забезпечує єдиний стандартизований інтерфейс.

% (Note) Колишній вставний розділ про «світовий досвід» інтегровано в основні глави.

\subsection*{1.4. Методологія Design Science Research}

Дослідження базується на методології \textbf{Design Science Research (DSR)} [16, 17]. На відміну від пояснювальних досліджень (explanatory research), DSR фокусується на створенні та оцінці інноваційних артефактів (моделей, методів, систем, прототипів), які вирішують практичні проблеми.

Hevner et al. сформулювали сім керівних принципів DSR \cite{Hevner2004DSR}:
\begin{enumerate}[leftmargin=*,itemsep=0.2em]
\item Design as an Artifact.
\item Problem Relevance.
\item Design Evaluation.
\item Research Contributions.
\item Research Rigor.
\item Design as a Search Process.
\item Communication of Research.
\end{enumerate}

Peffers et al. запропонували процесну модель DSR, яка включає шість етапів \cite{Peffers2007DSRMeth}.
\begin{longtable}{@{}p{0.14\textwidth}p{0.36\textwidth}p{0.44\textwidth}@{}}
\toprule
\textbf{Етап} & \textbf{Зміст} & \textbf{Реалізація в нашому дослідженні} \\
\midrule
1. Problem identification & Визначення проблеми та обґрунтування цінності рішення & Криза реактивної моделі бібліометричної звітності (розділ 2) \\
2. Define objectives & Формулювання вимог до рішення & Автоматизація, точність, інтеграція джерел, real-time аналітика \\
3. Design \& development & Проєктування та створення артефакту & Концепція MCP Agent, алгоритм Golden Record (розділ 4) \\
4. Demonstration & Демонстрація використання артефакту & Прототип ResearchIQ (розділ 5) \\
5. Evaluation & Оцінка ефективності & Кількісні та якісні метрики (розділ 5) \\
6. Communication & Публікація результатів & Дана стаття та подальші публікації \\
\bottomrule
\end{longtable}

Наше дослідження позиціонується на \textbf{Level 2} (nascent design theory) за класифікацією Gregor \& Hevner \cite{Gregor2013PositioningDSR}.

\textbf{Методи збору та аналізу даних:}

\textit{Кількісні методи:}
\begin{itemize}[leftmargin=*,itemsep=0.2em]
\item вимірювання часу виконання операцій (до/після впровадження);
\item аналіз якості даних (частка дублікатів, повнота заповнення);
\item порівняння охоплення джерел.
\end{itemize}

\textit{Якісні методи:}
\begin{itemize}[leftmargin=*,itemsep=0.2em]
\item аналіз user experience через спостереження за використанням системи;
\item експертна оцінка відповідності звітів вимогам МОН;
\item оцінка аналітичної цінності генерованих інсайтів.
\end{itemize}

\textbf{Технологічний стек апробації:}
\begin{itemize}[leftmargin=*,itemsep=0.2em]
\item Frontend: React/TypeScript для інтерактивного дашборду;
\item LLM: Google Gemini 3 Pro з підтримкою Function Calling;
\item Джерела даних: ORCID API v3, OpenAlex REST API;
\item Формати виводу: DOCX, XLSX, PDF з підтримкою шаблонів.
\end{itemize}